\section{GCD on Blackboard}
\label{sec:n-gcd-blackboard}
\source{AtCoder ABC125 C (Notion: Number Theory)}{Medium}

\problemstatement{Replace one of $N\le10^5$ integers ($\le10^9$) by any
integer in $[1,10^9]$ to maximise the gcd of all $N$ numbers.}

\begin{patternbox}
``Everything except element $i$'' for an associative operation $\Rightarrow$
\textbf{prefix and suffix} arrays.
\end{patternbox}

\subsection*{Approach and intuition}

If we replace $A_i$, the best new value is the gcd $g_i$ of all other numbers
itself (the gcd can never exceed $g_i$, and choosing $g_i$ attains it). So the
answer is $\max_i g_i$ with $g_i=\gcd(\text{pre}_{i-1},\text{suf}_{i+1})$,
using $\gcd(0,x)=x$ for the empty sides.

\subsection*{Algorithm}
\begin{algorithmic}[1]
\State $pre_0\gets0$, $pre_i\gets\gcd(pre_{i-1},A_i)$; $suf_{N+1}\gets0$, $suf_i\gets\gcd(suf_{i+1},A_i)$
\State \Return $\max_i\gcd(pre_{i-1},suf_{i+1})$
\end{algorithmic}

\dryrun{$[7,6,8]$: drop $7\to\gcd(6,8)=2$, drop $6\to1$, drop $8\to1$. Answer
$2$. $[12,15,18]$: drop $15\to6$. \checkmark}

\subsection*{C++17 implementation}
\cppfile{number-theory/12_gcd_on_blackboard.cpp}

\begin{complexitybox}
\textbf{Time:} $O(N\log V)$. \textbf{Space:} $O(N)$.
\end{complexitybox}

\begin{pitfallbox}
\begin{itemize}
  \item $N=2$: answer $\max(A_1,A_2)$ --- replace the smaller by the larger.
  \item Recomputing the gcd of the other $N-1$ elements for each $i$ is
        $O(N^2)$.
\end{itemize}
\end{pitfallbox}
